\chapter*{Resumen}
\addcontentsline{toc}{chapter}{Resumen}

El análisis espacial de fenómenos urbanos, como la criminalidad, requiere la asignación precisa de eventos georreferenciados sobre la infraestructura vial mediante procesos de \textit{Map Matching}. Sin embargo, los grafos urbanos extraídos de fuentes de datos abiertas como OpenStreetMap presentan una densidad geométrica masiva que incrementa significativamente la complejidad temporal y espacial de las consultas geográficas. 

Este trabajo de investigación propone y evalúa algoritmos de simplificación topológica de redes urbanas orientados a optimizar la eficiencia computacional en la asignación de datos delictivos. Para ello, se diseña e implementa una arquitectura modular que incorpora un \textit{pipeline} de evaluación estandarizado. Este marco permite cuantificar rigurosamente el impacto de la reducción de nodos y aristas a través de métricas objetivas, tales como el ratio de compresión topológica y la invarianza en la longitud de aristas. 

La solución computacional integra estructuras de indexación espacial avanzadas para acelerar las proyecciones geométricas en coordenadas planas. Los resultados obtenidos demuestran que es factible reducir drásticamente la huella de memoria del grafo y los tiempos de procesamiento de consultas espaciales, sin comprometer la precisión geométrica de la asignación ni alterar las invariantes topológicas necesarias para el análisis espacial urbano a gran escala.

\vspace{1cm}
\noindent \textbf{Palabras clave:} \textit{Map Matching}, Simplificación de Grafos, OpenStreetMap, Análisis Espacial, Estructuras de Datos Espaciales, Redes Urbanas.



\chapter*{Abstract}
\addcontentsline{toc}{chapter}{Abstract}

The spatial analysis of urban phenomena, such as crime, relies on the accurate assignment of georeferenced events to road infrastructure through Map Matching processes. However, urban graphs extracted from open data sources like OpenStreetMap exhibit a massive geometric density that significantly increases the temporal and spatial complexity of geographic queries.

This research proposes and evaluates topological simplification algorithms for urban networks, aimed at optimizing computational efficiency in the spatial assignment of crime data. To achieve this, a modular architecture incorporating a standardized evaluation pipeline is designed and implemented. This framework allows for the rigorous quantification of the impact of node and edge reduction through objective metrics, such as the topological compression ratio and edge length invariance.

The computational solution integrates advanced spatial indexing structures to accelerate geometric projections in planar coordinates. The results demonstrate that it is feasible to drastically reduce the memory footprint of the graph and the processing times for spatial queries, without compromising the geometric accuracy of the assignment or altering the fundamental topological invariants required for large-scale urban spatial analysis.

\vspace{1cm}
\noindent \textbf{Keywords:} Map Matching, Graph Simplification, OpenStreetMap, Spatial Analysis, Spatial Data Structures, Urban Networks.